\section{Demoralization}

\begin{rulebox}{General Rule}
Both players' armies have a demoralization level, representing the amount of punishment the army can take before beginning to disintegrate. During play, the elimination of friendly units and leaders and the enemy capture of artillery pieces and trains result in the accumulation of loss points. When the number of loss points accumulated by an army reaches that army's demoralization level, the army immediately becomes permanently demoralized. Demoralization severely limits the capabilities of an army's units.
\end{rulebox}

\ruleslabel{Procedure}

Each player individually keeps a running tally on a sheet of scratch paper of the number of loss points that his army has accumulated. As soon as an army reaches its demoralization level, the player commanding that army announces this fact. Thereafter, no further loss points are accumulated by either player---only one army can be demoralized per game.

\caseheading{13.1} The Royalist Army has a demoralization level of 100 loss points and the Allied Army has a demoralization level of 115 loss points.

\caseheading{13.2} Units become demoralized whenever their army reaches its demoralization level. All units of an army (regardless of the force to which they are attached) become demoralized whenever their army becomes demoralized.

\textsc{Exception:} The seven Newcastle foot regiments and the six Manchester horse regiments never become demoralized and are never affected by the provisions of this section.

\caseheading{13.3} Only one army per game can become demoralized.

As soon as one army becomes demoralized, all recording of loss points ceases for the rest of the game. Loss points are added to each army's total at the exact moment they are accumulated, and a player must immediately announce when his army has reached or exceeded its demoralization level. If both armies reach or exceed their demoralization levels at exactly the same time as a result of a combat, the inactive player's army is demoralized, and the active player's army is unaffected.

\caseheading{13.4} Once an army has become demoralized, it remains demoralized for the rest of the game. A demoralized army never ceases to suffer the effects of demoralization.

\caseheading{13.5} Loss points are accumulated by an army whenever its leaders or units are eliminated or its artillery pieces or train are captured.

\begin{tightitems}
  \item Whenever a unit is eliminated in combat (either due to an elimination result or because the unit suffered a disruption effect while already disrupted), a number of loss points equal to the combat strength printed on the unit's ordered (front) side is added to the accumulated loss-point total of its army.
  \item Whenever a leader other than Leven is eliminated for any reason, a number of loss points equal to five times the leader's command rating is added to the accumulated loss-point total of its army.
  \item Whenever an artillery piece is captured or recaptured, the player who gained control of it subtracts 5 loss points from his army's accumulated loss-point total, and his opponent adds 5 loss points to his army's accumulated loss-point total.
  \item Whenever a train hex is first entered by an enemy unit, the army to which the train belongs adds 20 loss points to its accumulated loss-point total, and the enemy army subtracts 5 loss points from its accumulated loss-point total.
\end{tightitems}

It is possible during a given game for both train hexes on the map to be entered, resulting in a loss-point addition to the army owning the train and a loss-point subtraction from the army entering the train hex. However, loss points are only added or subtracted as a result of the first entry of each individual train hex by an enemy unit. On the other hand, the loss-point additions and subtractions resulting from the capture or recapture of a particular artillery piece can occur any number of times during play.

\caseheading{13.6} Demoralized units are eligible for a rally check during their army's rally phase only if they are stacked in the same hex with a friendly leader of the same force. The foot units of Newcastle's force and the horse units of Manchester's force are always eligible for a rally check during their army's rally phase, since they never become demoralized.

\caseheading{13.7} All disrupted units of a demoralized army have their movement allowance doubled. Instead of moving 2 hexes per game turn, they move 4 hexes per game turn.

\caseheading{13.8} Demoralized units can leave the map.

Only leaders and demoralized units of a demoralized army can exit the map (the Allied leader Leven is an exception; see 12.9). Demoralized Royalist units and Royalist leaders can exit the map only via any hex on the northern map edge or any hex on the eastern map edge between hexes 1622 and 2622, inclusive. Demoralized Allied units and Allied leaders can exit the map only via any hex on the southern map edge.

Once a unit or leader exits the map, it can never re-enter the game. Disrupted units that exit the map cannot be rallied. They remain disrupted for the rest of the game and count for victory-point accumulation. Ordered units and leaders that exit the map have no further effect on the game; they do not count for victory-point accumulation. It costs both leaders and units (including disrupted units) one additional MP to exit the map.

\caseheading{13.9} Demoralized Heavy Horse units with a morale rating of less than 3 cannot charge. Only the Rupert, Byron, and RLG units of the Royalist Army can charge once that army is demoralized. No Allied Heavy Horse units other than those attached to Manchester's force (which do not become demoralized) can charge once the Allied Army has become demoralized.
